\documentclass[11pt, a4paper]{article}

% bfw-style.tex — gemeinsame Style-Definitionen für alle Handouts/Aufgabenhefte
% Voraussetzung: Kompilation mit xelatex (wegen fontspec)

\usepackage[ngerman]{babel}
\usepackage{geometry}
\geometry{a4paper, top=2.2cm, bottom=2.2cm, left=2.3cm, right=2.3cm}

\usepackage{fontspec}
% Fallback-Kette: Source Sans Pro -> Calibri -> Segoe UI -> Latin Modern Sans
% (Latin Modern ist immer mit MiKTeX vorhanden)
\IfFontExistsTF{Source Sans Pro}{%
  \setmainfont{Source Sans Pro}[Ligatures=TeX]%
}{\IfFontExistsTF{Calibri}{%
  \setmainfont{Calibri}[Ligatures=TeX]%
}{\IfFontExistsTF{Segoe UI}{%
  \setmainfont{Segoe UI}[Ligatures=TeX]%
}{%
  \setmainfont{Latin Modern Sans}[Ligatures=TeX]%
}}}
\IfFontExistsTF{Source Code Pro}{%
  \setmonofont{Source Code Pro}[Scale=0.92]%
}{\IfFontExistsTF{Consolas}{%
  \setmonofont{Consolas}[Scale=0.92]%
}{%
  \setmonofont{Latin Modern Mono}[Scale=0.92]%
}}

\usepackage{xcolor}
% Farbpalette: hell, freundlich, modern
\definecolor{bfwPrimary}{HTML}{0EA5A4}    % Petrol/Türkis - Primär
\definecolor{bfwPrimaryDark}{HTML}{0F766E} % dunkler Petrol für Headings
\definecolor{bfwAccent}{HTML}{F59E0B}     % Amber - Akzent
\definecolor{bfwSuccess}{HTML}{10B981}    % Grün - Erfolg / Lösung
\definecolor{bfwWarn}{HTML}{F97316}       % Orange - Achtung
\definecolor{bfwInk}{HTML}{0F172A}        % fast Schwarz
\definecolor{bfwInkSoft}{HTML}{475569}    % gedämpftes Grau
\definecolor{bfwBgSoft}{HTML}{F8FAFC}     % off-white
\definecolor{bfwBgTeal}{HTML}{ECFEFF}     % helles Türkis
\definecolor{bfwBgAmber}{HTML}{FEF3C7}    % helles Gelb
\definecolor{bfwBgRose}{HTML}{FFE4E6}     % helles Rosé für Eselsbrücken

\usepackage[colorlinks=true, linkcolor=bfwPrimaryDark, urlcolor=bfwPrimaryDark]{hyperref}
\usepackage{enumitem}
\setlist[itemize]{leftmargin=*, itemsep=2pt, topsep=2pt}
\setlist[enumerate]{leftmargin=*, itemsep=2pt, topsep=2pt}

\usepackage[most]{tcolorbox}
\usepackage{booktabs}
\usepackage{tikz}
\usetikzlibrary{decorations.pathmorphing, positioning, shapes.geometric, arrows.meta}

% Merkkästchen — wichtige Definition / Kernpunkt
\newtcolorbox{merksatz}[1][]{%
  enhanced, breakable,
  colback=bfwBgTeal,
  colframe=bfwPrimary,
  boxrule=0.6pt,
  arc=4pt,
  left=10pt, right=10pt, top=8pt, bottom=8pt,
  fonttitle=\bfseries\color{bfwPrimaryDark},
  title={\faLightbulb\ Merksatz},
  #1
}

% Eselsbrücke — Mnemoniks
\newtcolorbox{eselsbruecke}[1][]{%
  enhanced, breakable,
  colback=bfwBgRose,
  colframe=bfwAccent,
  boxrule=0.6pt,
  arc=4pt,
  left=10pt, right=10pt, top=8pt, bottom=8pt,
  fonttitle=\bfseries\color{bfwWarn},
  title={Eselsbrücke},
  #1
}

% Beispiel-Box
\newtcolorbox{beispiel}[1][]{%
  enhanced, breakable,
  colback=bfwBgSoft,
  colframe=bfwInkSoft,
  boxrule=0.4pt,
  arc=3pt,
  left=10pt, right=10pt, top=8pt, bottom=8pt,
  fonttitle=\bfseries\color{bfwInk},
  title={Beispiel},
  #1
}

% Lösungs-Box (für Aufgabenhefte)
\newtcolorbox{loesung}[1][]{%
  enhanced, breakable,
  colback=bfwBgAmber!40,
  colframe=bfwSuccess,
  boxrule=0.6pt,
  arc=3pt,
  left=10pt, right=10pt, top=8pt, bottom=8pt,
  fonttitle=\bfseries\color{bfwSuccess!70!black},
  title={Lösung},
  #1
}

% Glossar-Eintrag
\newcommand{\glossareintrag}[2]{%
  \par\noindent\textbf{\color{bfwPrimaryDark}#1} \enspace #2\par\smallskip
}

% Section-Stil — etwas farbiger als Standard
\usepackage{titlesec}
\titleformat{\section}
  {\Large\bfseries\color{bfwPrimaryDark}}
  {\thesection}{0.6em}{}
\titleformat{\subsection}
  {\large\bfseries\color{bfwPrimary}}
  {\thesubsection}{0.5em}{}
\titleformat{\subsubsection}
  {\normalsize\bfseries\color{bfwInk}}
  {\thesubsubsection}{0.4em}{}

% TikZ-Stil "skizzenhaft" — etwas weicher als Rough.js, aber stabil
\tikzset{
  roughbox/.style = {
    draw=bfwPrimaryDark, thick, fill=bfwBgTeal,
    rounded corners=4pt, inner sep=6pt
  },
  roughaccent/.style = {
    draw=bfwAccent, thick, fill=bfwBgAmber,
    rounded corners=4pt, inner sep=6pt
  },
  roughline/.style = {
    draw=bfwInkSoft, thick, -{Stealth[length=2.5mm]}
  }
}

% Bullet-Symbole (bewusst kein FontAwesome — vermeidet Font-Installations-Probleme)
\newcommand{\faLightbulb}{$\bullet$}
\newcommand{\faBookmark}{$\bullet$}

% Header/Footer
\usepackage{fancyhdr}
\pagestyle{fancy}
\fancyhf{}
\renewcommand{\headrulewidth}{0pt}
\fancyfoot[L]{\small\color{bfwInkSoft}\thepage}
\fancyfoot[R]{\small\color{bfwInkSoft} BFW M-064 Beschaffung im Büromanagement}


\title{\vspace*{-1.5cm}\Huge\bfseries\color{bfwPrimaryDark}Aufgabenheft Tag~12\\[0.4em]
  \large\color{bfwInkSoft}\textnormal{Dokumentenmanagement: Aufbewahrung \& Fristen --- Übungen im IHK-Klausurformat}}
\author{\textsf{Büroprozesse gestalten und Arbeitsvorgänge organisieren \textperiodcentered{} M-067}\\
\textsf{\small\copyright\ Can Siebert\,\textperiodcentered\,2026}}
\date{Selbstlernmaterial \textperiodcentered{} 14.07.2026}

\begin{document}
\maketitle
\thispagestyle{empty}

\begin{merksatz}[title={\bfseries So nutzen Sie dieses Aufgabenheft}]
Bearbeiten Sie alle Aufgaben zunächst \emph{ohne} in die Lösungen zu schauen. Schreiben
Sie Ihre Antworten in vollständigen Sätzen --- die IHK-Prüfung verlangt formulierte
Begründungen, nicht bloße Stichworte. Beachten Sie die \emph{Operatoren} (nennen,
erläutern, beurteilen, berechnen, begründen) --- sie geben den geforderten
Antwort-Tiefegrad vor. Erst danach öffnen Sie den Lösungsteil ab Seite~\pageref{loesungen}
und vergleichen.
\end{merksatz}

\tableofcontents
\vspace{1em}
\hrule

\section{Aufgaben}

\subsection*{Aufgabe~1 --- Bedeutung der Schriftgutaufbewahrung (10 Punkte)}

\noindent\textbf{\color{bfwAccent}YouTube-Vorbereitung:}\enspace \href{https://www.youtube.com/watch?v=zx9vxdezRyc}{Schriftgutverwaltung --- Grundlagen} (\url{https://www.youtube.com/watch?v=zx9vxdezRyc})

\medskip
In der \emph{Duisdorfer BüroKonzept KG} fallen täglich zahlreiche Schriftstücke an. Die
Auszubildende fragt, warum man die Unterlagen nicht einfach nach Erledigung entsorgt.

\begin{enumerate}[label=(\alph*)]
  \item \textbf{Erläutern} Sie, warum ein Unternehmen Schriftstücke geordnet aufbewahren muss. \hfill (3~P.)
  \item \textbf{Nennen} Sie je einen betrieblichen, rechtlichen und steuerlichen Grund der Aufbewahrung. \hfill (3~P.)
  \item \textbf{Erklären} Sie, warum die Registratur als „Gedächtnis des Unternehmens” bezeichnet wird. \hfill (2~P.)
  \item \textbf{Beurteilen} Sie die Aussage: \emph{„Belege kann man nach der Buchung gleich vernichten.”} \hfill (2~P.)
\end{enumerate}

\subsection*{Aufgabe~2 --- GoBD und digitale Aufbewahrung (10 Punkte)}

\noindent\textbf{\color{bfwAccent}YouTube-Vorbereitung:}\enspace \href{https://www.youtube.com/watch?v=r1kpfyHPcvI}{GoBD einfach erklärt} (\url{https://www.youtube.com/watch?v=r1kpfyHPcvI})

\medskip
Die KG möchte eingehende Rechnungen künftig nur noch digital archivieren.

\begin{enumerate}[label=(\alph*)]
  \item \textbf{Nennen} Sie, wofür die Abkürzung \emph{GoBD} steht. \hfill (2~P.)
  \item \textbf{Erläutern} Sie zwei Anforderungen nach HGB § 257 Abs.~3 an die digitale Aufbewahrung. \hfill (4~P.)
  \item \textbf{Begründen} Sie, warum die \emph{Unveränderbarkeit} der Belege so wichtig ist. \hfill (2~P.)
  \item \textbf{Beurteilen} Sie, ob ein eingescanntes PDF einer Papierrechnung gleichwertig sein kann. \hfill (2~P.)
\end{enumerate}

\subsection*{Aufgabe~3 --- Wertigkeitsstufen (12 Punkte)}

\noindent\textbf{\color{bfwAccent}YouTube-Vorbereitung:}\enspace \href{https://www.youtube.com/watch?v=wsj-jRTouXY}{Wertigkeitsstufen von Schriftgut} (\url{https://www.youtube.com/watch?v=wsj-jRTouXY})

\medskip
Schriftstücke werden nach ihrer Bedeutung vier Wertigkeitsstufen zugeordnet.

\begin{enumerate}[label=(\alph*)]
  \item \textbf{Nennen} Sie die vier Wertigkeitsstufen. \hfill (2~P.)
  \item \textbf{Erläutern} Sie kurz jede Stufe mit einem Beispiel. \hfill (4~P.)
  \item \textbf{Ordnen} Sie zu: (1)~Tageszeitung, (2)~verlangtes Angebot, (3)~Ausgangsrechnung, (4)~Patentunterlage. \hfill (4~P.)
  \item \textbf{Erläutern} Sie den Unterschied zwischen \emph{Prüfwert} und \emph{Gesetzeswert}. \hfill (2~P.)
\end{enumerate}

\subsection*{Aufgabe~4 --- Aufbewahrungsfristen berechnen (16 Punkte)}

\noindent\textbf{\color{bfwAccent}YouTube-Vorbereitung:}\enspace \href{https://www.youtube.com/watch?v=XHUau_xKmaY}{Aufbewahrungsfristen HGB \& AO} (\url{https://www.youtube.com/watch?v=XHUau_xKmaY})

\medskip
Beachten Sie: Buchungsbelege sind seit dem Vierten Bürokratieentlastungsgesetz (BEG~IV,
in Kraft seit 2025) nur noch 8 Jahre aufzubewahren; Bücher und Jahresabschlüsse weiterhin
10 Jahre, Handelsbriefe 6 Jahre.

\begin{enumerate}[label=(\alph*)]
  \item \textbf{Nennen} Sie die Aufbewahrungsfristen für Jahresabschlüsse, Buchungsbelege und Handelsbriefe. \hfill (3~P.)
  \item \textbf{Erläutern} Sie, wann die Aufbewahrungsfrist beginnt. \hfill (3~P.)
  \item \textbf{Berechnen} Sie das Fristende für eine Ausgangsrechnung (Buchungsbeleg), die am 05.\,06.\,2025 erstellt wurde. \hfill (3~P.)
  \item \textbf{Berechnen} Sie das Fristende für einen im Jahr 2024 empfangenen Handelsbrief. \hfill (3~P.)
  \item \textbf{Berechnen} Sie das Fristende für den Jahresabschluss des Geschäftsjahres 2025 (festgestellt 2026). \hfill (2~P.)
  \item \textbf{Beurteilen} Sie, welchen Vorteil die Verkürzung auf 8 Jahre für Unternehmen bringt. \hfill (2~P.)
\end{enumerate}

\subsection*{Aufgabe~5 --- Schriftgutkatalog (12 Punkte)}

\noindent\textbf{\color{bfwAccent}YouTube-Vorbereitung:}\enspace \href{https://www.youtube.com/watch?v=JnbejlXR55Q}{Schriftgutkatalog \& Aktenplan} (\url{https://www.youtube.com/watch?v=JnbejlXR55Q})

\medskip
Die \emph{Duisdorfer BüroKonzept KG} legt für die Abteilung Vertrieb einen Schriftgutkatalog an.

\begin{enumerate}[label=(\alph*)]
  \item \textbf{Erklären} Sie den Zweck eines Schriftgutkatalogs. \hfill (3~P.)
  \item \textbf{Nennen} Sie die Angaben, die ein Schriftgutkatalog je Schriftstück enthält. \hfill (3~P.)
  \item \textbf{Ordnen} Sie folgenden Schriftstücken Wertigkeitsstufe und Frist zu: Ausgangsrechnungen, verlangte Angebote, Preislisten, Gründungsurkunde. \hfill (4~P.)
  \item \textbf{Beurteilen} Sie, warum unnötig langes Aufbewahren vermieden werden sollte. \hfill (2~P.)
\end{enumerate}

\subsection*{Aufgabe~6 --- Ablageorte und Registraturarten (10 Punkte)}

\noindent\textbf{\color{bfwAccent}YouTube-Vorbereitung:}\enspace \href{https://www.youtube.com/watch?v=ldee6qhCJe8}{Zentrale \& dezentrale Ablage} (\url{https://www.youtube.com/watch?v=ldee6qhCJe8})

\medskip
\begin{enumerate}[label=(\alph*)]
  \item \textbf{Nennen} Sie den Grundsatz, nach dem der Ablageort gewählt wird. \hfill (2~P.)
  \item \textbf{Erläutern} Sie kurz Arbeitsplatzablage, Abteilungsablage und Zentralregistratur/Archiv. \hfill (3~P.)
  \item \textbf{Ordnen} Sie zu \textbf{und begründen} Sie: \hfill (3~P.)
    \begin{enumerate}[label=\arabic*., leftmargin=*]
      \item Eine laufende Kundenakte mit täglichem Zugriff.
      \item Unterlagen, die mehrere Mitarbeiter einer Abteilung gemeinsam nutzen.
      \item Buchungsbelege aus 2019, die nur noch wegen der gesetzlichen Frist aufbewahrt werden.
    \end{enumerate}
  \item \textbf{Unterscheiden} Sie zentrale und dezentrale Ablage anhand je eines Vorteils. \hfill (2~P.)
\end{enumerate}

\vspace{1em}
\noindent\textbf{Punkte gesamt: 70}

\newpage

\section{Lösungen}\label{loesungen}

\subsection*{Lösung Aufgabe~1}
\begin{loesung}
\textbf{(a)} Damit schnell und unkompliziert auf benötigte Unterlagen zurückgegriffen
werden kann, ergibt sich die innerbetriebliche Notwendigkeit einer geordneten
Aufbewahrung. Zusätzlich bestehen gesetzliche Vorschriften, die die Aufbewahrung von
Schriftstücken zur Dokumentation und Beweissicherung regeln.

\textbf{(b)} Betrieblich: schneller Zugriff auf Vorgänge und Korrespondenz. Rechtlich:
Dokumentation und Beweissicherung der Geschäftsvorfälle nach HGB. Steuerlich:
Buchführungsunterlagen müssen für Prüf- und Beweiszwecke (Finanzamt) verfügbar sein (AO).

\textbf{(c)} Weil in der Registratur das gesamte „Wissen” über vergangene
Geschäftsvorfälle gespeichert ist und jederzeit abgerufen werden kann --- sie ist das
Gedächtnis des Unternehmens.

\textbf{(d)} Die Aussage ist falsch. Buchungsbelege haben Gesetzeswert und unterliegen
gesetzlichen Aufbewahrungsfristen (Buchungsbelege 8 Jahre). Eine sofortige Vernichtung
wäre ein Verstoß gegen HGB § 257 / AO § 147.
\end{loesung}

\subsection*{Lösung Aufgabe~2}
\begin{loesung}
\textbf{(a)} GoBD = Grundsätze zur ordnungsmäßigen Führung und Aufbewahrung von Büchern,
Aufzeichnungen und Unterlagen in elektronischer Form sowie zum Datenzugriff.

\textbf{(b)} Zwei Anforderungen: (1)~Die Wiedergabe muss mit empfangenen Handelsbriefen
und Buchungsbelegen \emph{bildlich} und mit den übrigen Unterlagen \emph{inhaltlich}
übereinstimmen, wenn sie lesbar gemacht wird. (2)~Die Unterlagen müssen während der
gesamten Aufbewahrungsfrist verfügbar und jederzeit innerhalb angemessener Frist lesbar
gemacht werden können.

\textbf{(c)} Nur unveränderbare Belege haben Beweiskraft. Könnte man Belege nachträglich
ändern, wäre eine lückenlose und manipulationssichere Nachvollziehbarkeit der
Geschäftsvorfälle nicht mehr gewährleistet --- das Finanzamt würde die Buchführung nicht
anerkennen.

\textbf{(d)} Ja, ein digitales Abbild kann gleichwertig sein, wenn es den GoBD entspricht:
bildliche Übereinstimmung mit dem Original, jederzeitige Lesbarkeit und Verfügbarkeit
während der gesamten Frist sowie revisionssichere, unveränderbare Speicherung.
\end{loesung}

\subsection*{Lösung Aufgabe~3}
\begin{loesung}
\textbf{(a)} Tageswert, Prüfwert, Gesetzeswert, Dauerwert.

\textbf{(b)}
\begin{itemize}[itemsep=0pt]
  \item \emph{Tageswert}: kein bleibender Wert, nach Kenntnisnahme vernichtbar (z.\,B.\ Tageszeitung).
  \item \emph{Prüfwert}: nur begrenzt bedeutsam, nach Prüfung entscheiden (z.\,B.\ Preisliste).
  \item \emph{Gesetzeswert}: gesetzliche Aufbewahrungsfrist nach HGB/AO (z.\,B.\ Rechnung).
  \item \emph{Dauerwert}: solange das Unternehmen besteht (z.\,B.\ Gründungsunterlage).
\end{itemize}

\textbf{(c)} (1)~Tageswert, (2)~Prüfwert, (3)~Gesetzeswert, (4)~Dauerwert.

\textbf{(d)} \emph{Prüfwert}: das Unternehmen entscheidet nach einer Prüfung selbst über
Vernichtung oder Aufbewahrung; es gibt keine gesetzliche Pflicht. \emph{Gesetzeswert}:
der Gesetzgeber schreibt eine verbindliche Aufbewahrungsfrist vor.
\end{loesung}

\subsection*{Lösung Aufgabe~4}
\begin{loesung}
\textbf{(a)} Jahresabschlüsse: 10 Jahre; Buchungsbelege: 8 Jahre (seit BEG~IV);
Handelsbriefe: 6 Jahre.

\textbf{(b)} Die Frist beginnt mit dem \emph{Schluss des Kalenderjahres}, in dem das
Schriftstück entstanden bzw.\ der Handelsbrief empfangen/abgesandt wurde --- also stets
zum 31.\,12. Das Fristende liegt deshalb immer auf einem 31.\,Dezember.

\textbf{(c)} Buchungsbeleg vom 05.\,06.\,2025: Fristbeginn 31.\,12.\,2025, 8 Jahre $\Rightarrow$
Fristende \textbf{31.\,12.\,2033}.

\textbf{(d)} Handelsbrief, empfangen 2024: Fristbeginn 31.\,12.\,2024, 6 Jahre $\Rightarrow$
Fristende \textbf{31.\,12.\,2030}.

\textbf{(e)} Jahresabschluss des Geschäftsjahres 2025 (festgestellt 2026): Fristbeginn
31.\,12.\,2026, 10 Jahre $\Rightarrow$ Fristende \textbf{31.\,12.\,2036}.

\textbf{(f)} Vorteil: geringerer Lager-, Archivierungs- und Verwaltungsaufwand sowie
schnellere Vernichtung nicht mehr benötigter Belege --- das spart Platz und Kosten
(Bürokratieabbau).
\end{loesung}

\subsection*{Lösung Aufgabe~5}
\begin{loesung}
\textbf{(a)} Der Schriftgutkatalog vermeidet Doppel- und Mehrfachablagen sowie unnötig
langes Aufbewahren. Er schafft Überblick über alle in einem Büro vorkommenden
Schriftstücke.

\textbf{(b)} Je Schriftstück: Wertigkeitsstufe, Aufbewahrungsfrist (gesetzlich/betrieblich),
Ablageort sowie ggf.\ Fristende und Bemerkungen.

\textbf{(c)}
\begin{itemize}[itemsep=0pt]
  \item Ausgangsrechnungen: Gesetzeswert, Frist (Buchungsbeleg) 8 Jahre.
  \item Verlangte Angebote: Prüfwert, keine gesetzliche Frist (betrieblich, Gültigkeit beachten).
  \item Preislisten: Prüfwert, keine gesetzliche Frist (Gültigkeit beachten).
  \item Gründungsurkunde: Dauerwert, dauerhaft aufzubewahren.
\end{itemize}

\textbf{(d)} Unnötig langes Aufbewahren bindet Platz und Kosten und erschwert das
Auffinden aktueller Unterlagen. Nach Ablauf der Frist können Unterlagen ohne Gesetzes-
oder Dauerwert vernichtet werden.
\end{loesung}

\subsection*{Lösung Aufgabe~6}
\begin{loesung}
\textbf{(a)} Grundsatz: Je häufiger ein Mitarbeiter auf die Schriftstücke zurückgreifen
muss, umso näher an seinem Arbeitsplatz sollten sie aufbewahrt werden.

\textbf{(b)}
\begin{itemize}[itemsep=0pt]
  \item \emph{Arbeitsplatzablage}: direkt am Platz, für häufig/schnell benötigte Unterlagen (dezentral).
  \item \emph{Abteilungsablage}: gemeinsam genutzt, für alle erreichbar und nachvollziehbar (Aktenplan).
  \item \emph{Zentralregistratur/Archiv}: Unterlagen verschiedener Bereiche; entlastet die Abteilungen, bewahrt auch Dauerwert auf.
\end{itemize}

\textbf{(c)} Zuordnung mit Begründung:
\begin{enumerate}[label=\arabic*., itemsep=0pt]
  \item Arbeitsplatzablage (Tagesregistratur) --- täglicher Zugriff, daher nah am Platz.
  \item Abteilungsablage --- gemeinsame Nutzung durch mehrere Mitarbeiter.
  \item Zentralregistratur/Archiv (Altregistratur) --- nur noch wegen der Aufbewahrungspflicht, seltener Zugriff.
\end{enumerate}

\textbf{(d)} \emph{Dezentrale} Ablage: schneller Zugriff am Arbeitsplatz/in der Abteilung.
\emph{Zentrale} Ablage: einheitliche Ordnung und Entlastung der Arbeitsplätze von selten
benötigten Unterlagen.
\end{loesung}

\end{document}
